\section{问题形式化}
\label{subsec:formalisation_mtdmdp}

\subsection{定义}

在预备知识中，我们介绍了两种转移概率混合模型：\Cref{def:mixture_transition_distribution} 的 \gls{mtd} 和 \Cref{def:multimatrix_mixture_transistion_distribution} 的 \gls{mtdg}。基于这些模型，我们将构建延迟过程，其中转移被定义为转移概率的混合。定义此类模型有多种方式，下面将探索其中四种。我们用术语 \gls{mtdmdp} 统称这些过程。现在介绍第一个模型。

\begin{defi}[\Gls{ismmdp}]
    给定一个转移为 $p$ 的 \gls{mdp} $\markovdecisionprocess$ 和延迟 $\delay$，\gls{ismmdp} 转移过程定义如下，
    \begin{align*}
        \probability(S_{t+1}=s_{t+1}\vert& S_{t}=s_{t},A_{t}=a_{t},\dots,A_{t-\delaymax}=a_{t-\delaymax})
        \\
        &=\sum_{g=0}^\delaymax \lambda_g \probability(S_{t+1}=s_{t+1}\vert S_{t}=s_{t},A_{t}=a_{t-g})
        \\
        &= \sum_{g=0}^\delaymax \lambda_g p(s_{t+1}\vert s_{t},a_{t-g}).
    \end{align*}
    其奖励为，
    \begin{align*}
        R(s_{t},a_{t},\dots,a_{t-\delaymax}) = \sum_{g=0}^\delaymax \lambda_g \rewardfunction(s_{t},a_{t-g}).
    \end{align*}
\end{defi}


\begin{remark}
    经典的常值 $\delay$-延迟 \gls{dmdp} 属于此类，只需令 $\lambda_0=\dots=\lambda_{\delay-1}=0$ 且 $\lambda_\delay=1$。
\end{remark}

第一个定义基于如下假设：无论延迟动作是哪个，它都将应用于当前状态。因此，动作作用于\emph{瞬时（instantaneous）}状态而非\emph{过去（past）}状态。这对应了名称中的字母"I"。第二个重要的字母是"S"，因为该模型基于常规的 \gls{mtd} 模型构建，该模型考虑\emph{单个（single）}转移矩阵。现在考虑基于 \gls{mtd} 的另一个模型，但使用过去状态来应用动作，因此名称中含"P"。

\begin{defi}[\Gls{psmmdp}]
    给定一个转移为 $p$ 的 \gls{mdp} $\markovdecisionprocess$ 和延迟 $\delay$，\gls{psmmdp} 转移过程定义如下，
    \begin{align*}
        \probability(S_{t+1}=s_{t+1}\vert& S_{t}=s_{t},A_{t}=a_{t},\dots,S_{t-\delaymax}=s_{t-\delaymax},A_{t-\delay}=a_{t-\delaymax})
        \\
        &=\sum_{g=0}^\delaymax \lambda_g \probability(S_{t+1}=s_t\vert S_{t}=s_{t-g},A_{t}=a_{t-g})
        \\
        &= \sum_{g=0}^\delaymax \lambda_g \transitionfunction(s_t\vert s_{t-g},a_{t-g}).
    \end{align*}
    其奖励为，
    \begin{align*}
        R(s_{t},a_{t},\dots,s_{t-\delaymax},a_{t-\delaymax}) = \sum_{g=0}^\delaymax \lambda_g \rewardfunction(s_{t-g},a_{t-g}).
    \end{align*}
\end{defi}

注意转移和奖励如何依赖于过去状态 $s_{t-g}$。现在定义最后两个模型，它们使用 \gls{mtdg} 模型，增加了更多自由度。由于它们依赖于\emph{多个（multiple）}转移矩阵，我们在名称中使用字母"M"。

\begin{defi}[\Gls{immmdp}]
    对于延迟 $\delaymax$ 以及各自转移为 $(p_g)_{g\in[\![1,\delaymax]\!]}$ 的 \glspl{mdp} $(\markovdecisionprocess_g)_{g\in[\![1,\delaymax]\!]}$，\gls{immmdp} 转移过程定义如下，
    \begin{align*}
        \probability(S_{t+1}=s_{t+1}\vert& S_{t}=s_{t},A_{t}=a_{t},\dots,A_{t-\delaymax}=a_{t-\delaymax})
        \\
        &=\sum_{g=0}^\delaymax \lambda_g \probability(S_t=s_t\vert S_{t}=s_{t},A_{t-g}=a_{t-g})
        \\
        &= \sum_{g=0}^\delaymax \lambda_g \transitionfunction_g(s_{t+1}\vert s_{t},a_{t-g}).
    \end{align*}
    其奖励为，
    \begin{align*}
        R(s_{t},a_{t},\dots,a_{t-\delaymax}) = \sum_{g=0}^\delaymax \lambda_g \rewardfunction(s_{t},a_{t-g}).
    \end{align*}
\end{defi}

\newpage
\begin{defi}[\Gls{pmmmdp}]
\label{def:pmmmdp}
    对于延迟 $\delaymax$ 以及各自转移为 $(p_g)_{g\in[\![1,\delaymax]\!]}$ 的 \glspl{mdp} $(\markovdecisionprocess_g)_{g\in[\![1,\delaymax]\!]}$，\gls{pmmmdp} 转移过程定义如下，
    \begin{align*}
        \probability(S_{t+1}=s_{t+1}\vert& S_{t}=s_{t},A_{t}=a_{t},\dots,S_{t-\delaymax}=s_{t-\delaymax},A_{t-\delay}=a_{t-\delaymax})
        \\
        &=\sum_{g=0}^\delaymax \lambda_g \probability(S_{t+1}=s_{t+1}\vert S_{t-g}=s_{t-g},A_{t-g}=a_{t-g})
        \\
        &= \sum_{g=0}^\delaymax \lambda_g \transitionfunction_g(s_{t+1}\vert s_{t-g},a_{t-g}).
    \end{align*}
    其奖励为，
    \begin{align*}
        R(s_{t},a_{t},\dots,s_{t-\delaymax},a_{t-\delaymax}) = \sum_{g=0}^\delaymax \lambda_g \rewardfunction(s_{t-g},a_{t-g}).
    \end{align*}
\end{defi}

注意初始状态分布尚未定义。一种可能性——本章采用的方式——是将过程初始化为常值 $\delaymax$-延迟 \gls{dmdp}：前 $\delaymax$ 个动作从动作空间中均匀随机采样。\Cref{tab:mtdmdp_properties} 总结了每个过程的性质。

\begin{table}[t]
\centering
    \begin{tabular}{cc|c|c|}
        \cline{3-4}
        &&\multicolumn{2}{c|}{\textbf{转移使用的状态}}\\
        \cline{3-4}
        & & 过去 & 瞬时 \\
        \hline
        \multicolumn{1}{|c|}{\multirow{2}{*}{\makecell{\textbf{转移}\\\textbf{矩阵}}}}&单一&\gls{psmmdp}&\gls{ismmdp}\\
        \cline{2-4}
        \multicolumn{1}{|c|}{}&多个&\gls{pmmmdp}&\gls{immmdp}\\
        \hline
    \end{tabular}
    \caption{从 \gls{mtd} 和 \gls{mtdg} 模型应用于 \glspl{mdp} 导出的四种过程的主要性质。}
    \label{tab:mtdmdp_properties}
 \end{table}


\subsection{目标与假设}
注意我们并不假设智能体知道每步哪个动作已被应用于环境，即我们处于匿名框架中。关于 \gls{mtdmdp} 框架下的目标，我们将研究无限时域下的期望折扣回报或平均奖励准则。特别地，我们将关注延迟向量 $\boldsymbol\lambda=(\lambda_0,\dots,\lambda_\delaymax)$ 的选择对智能体性能的影响。关于信息结构（见 \Cref{subsec:control_theory}），与经典常值延迟 $\delay$ 类似，我们假设智能体在时刻 $t$ 可以访问历史 $h_t = (h^s_t,h^r_t,h^a_t)$，其中 $h^s_t = (s_i)_{0\le i\le t-\delay}$ 为状态历史，$h^r_t = (r_i)_{0\le i\le t-\delay}$ 为奖励历史，$h^a_t = (r_i)_{0\le i\le t}$ 为动作历史。

在理论分析中，我们假设状态空间和动作空间有限，以简化问题并利用理论 \gls{rl} 的结果。


\subsection{记号}
我们将采用 \Cref{subsec:def_state_distrib} 中定义的面向期望折扣回报和平均奖励目标的状态分布。由于某些结果可通过类似计算同时适用于两者，为简洁起见，我们可能在这些情况下省略下标"AVG"或"$\gamma$"。此外，我们使用记号 $\augmentedstatespace=\statespace\times\actionspace^{\delaymax}$ 表示 \gls{ismmdp} 和 \gls{immmdp} 的增广状态空间，类似于常值延迟 \gls{dmdp}。对于 \gls{psmmdp} 和 \gls{pmmmdp}，我们将增广状态空间定义为 $\orderdstatespace=\statespace^{\delaymax}\times\actionspace^{\delaymax}$。由于与常值延迟 \gls{dmdp} 的相似性，容易证明 \gls{ismmdp} 和 \gls{immmdp} 中的策略属于集合 $\delayedpolicyspace$~（参见 \Cref{subsubsec:theory_augmented_space}）。\gls{psmmdp} 和 \gls{pmmmdp} 在其转移模型中考虑过去状态，因此它们不会向增广状态添加比 \gls{ismmdp} 和 \gls{immmdp} 的增广状态中已有的更多最新信息。这意味着前者的策略从后者的角度来看是历史依赖的，也属于 $\delayedpolicyspace$。我们将它们的策略集记为 $\higherorderpolicyspace \subset \delayedpolicyspace$。为避免混淆，我们将 \gls{psmmdp} 和 \gls{pmmmdp} 的策略记作 $\orderdpolicy$。

最后，定义 \gls{mtdmdp} 中状态和动作分布的记号。考虑 $s\in\statespace$、$a\in\actionspace$、$x\in\augmentedstatespace$、$\widebar x\in\orderdstatespace$、$\pi\in\Pi$ 以及 $\delayedpolicy,\orderdpolicy\in\delayedpolicyspace$。那么，无延迟策略 $\pi$ 的分布写作，
\begin{enumerate}
    \item $\discountedstateoccupancydistribution[][\pi](s)$ 为 \gls{mdp} 中策略 $\pi$ 下的状态分布；
    \item $\discountedstateoccupancydistribution[][\pi](s,a)$ 为 \gls{mdp} 中策略 $\pi$ 下的状态-动作分布；
\end{enumerate}
对于 \gls{ismmdp} 和 \gls{immmdp}，分布为，\footnote{可以看出，我们选择不在记号中显式指明输入空间，以保持可读性。输入空间由分布所评估的量即可明确。}
\begin{enumerate}
\setcounter{enumi}{2}
    \item $\delayeddiscountedstateoccupancydistribution[][\delayedpolicy](x)$ 为 \gls{ismmdp} 和 \gls{immmdp} 中策略 $\delayedpolicy$ 下在 $\augmentedstatespace$ 上的分布；
    \item $\delayeddiscountedstateoccupancydistribution[][\delayedpolicy](x,a)$ 为 \gls{ismmdp} 和 \gls{immmdp} 中策略 $\delayedpolicy$ 下在 $\augmentedstatespace\times\actionspace$ 上的分布。
    \item $\delayeddiscountedstateoccupancydistribution[][\delayedpolicy](s)$ 为 \gls{ismmdp} 和 \gls{immmdp} 中策略 $\delayedpolicy$ 下在 $\statespace$ 上的分布；
    \item $\delayeddiscountedstateoccupancydistribution[][\delayedpolicy](s,a)$ 为 \gls{ismmdp} 和 \gls{immmdp} 中策略 $\delayedpolicy$ 下在 $\statespace\times\actionspace$ 上的分布。
\end{enumerate}
类似地，对于 \gls{psmmdp} 和 \gls{pmmmdp}，分布为，
\begin{enumerate}
\setcounter{enumi}{6}
    \item $\delayeddiscountedstateoccupancydistribution[][\orderdpolicy](\widebar x)$ 为 \gls{psmmdp} 和 \gls{pmmmdp} 中策略 $\orderdpolicy$ 下在 $\orderdstatespace$ 上的分布；
    \item $\delayeddiscountedstateoccupancydistribution[][\orderdpolicy](\widebar x,a)$ 为 \gls{psmmdp} 和 \gls{pmmmdp} 中策略 $\orderdpolicy$ 下在 $\orderdstatespace\times\actionspace$ 上的分布。
    \item $\delayeddiscountedstateoccupancydistribution[][\orderdpolicy](s)$ 为 \gls{psmmdp} 和 \gls{pmmmdp} 中策略 $\orderdpolicy$ 下在 $\statespace$ 上的分布；
    \item $\delayeddiscountedstateoccupancydistribution[][\orderdpolicy](s,a)$ 为 \gls{psmmdp} 和 \gls{pmmmdp} 中策略 $\orderdpolicy$ 下在 $\statespace\times\actionspace$ 上的分布。
\end{enumerate}